\documentclass[a4paper,14pt]{extarticle}

% ==================== KHAI BÁO GÓI ====================
\usepackage{fontspec}
\setmainfont{TeX Gyre Termes}

\usepackage[vietnamese]{babel}
\usepackage[
    a4paper,
    left=1.5cm,
    right=1.5cm,
    top=1.8cm,
    bottom=1.7cm,
    headheight=14pt,
    headsep=0.35cm,
    footskip=0.7cm
]{geometry}

\usepackage{fancyhdr}
\usepackage{tabularx}
\usepackage{array}
\usepackage{booktabs}
\usepackage{enumitem}
\usepackage{microtype}
\usepackage{ragged2e}
\usepackage{setspace}

% ==================== ĐỊNH DẠNG CHUNG ====================
\setlength{\parindent}{0pt}
\setlength{\parskip}{3pt}
\setlength{\tabcolsep}{3pt}
\renewcommand{\arraystretch}{1.08}

\newcolumntype{Y}{>{\RaggedRight\arraybackslash}X}
\newcolumntype{P}[1]{>{\RaggedRight\arraybackslash}p{#1}}

\setlist[itemize]{
    leftmargin=1.1cm,
    itemsep=2pt,
    topsep=2pt
}

% ==================== ĐẦU TRANG / CHÂN TRANG ====================
\pagestyle{fancy}
\fancyhf{}

\fancyhead[L]{\small Problem Analysis -- Smart Money Bot (bản hoàn chỉnh)}
\fancyfoot[R]{\small Page \thepage\ of 10}

\renewcommand{\headrulewidth}{0pt}
\renewcommand{\footrulewidth}{0pt}

% ==================== LỆNH ĐỊNH DẠNG BẢNG ====================
\newcommand{\TableText}{\footnotesize}
\newcommand{\TableStart}{%
    \begingroup
    \TableText
    \renewcommand{\arraystretch}{1.05}
}
\newcommand{\TableEnd}{\endgroup}

\begin{document}

% ============================================================
% TRANG 1
% ============================================================

\begin{center}
    {\Large\bfseries BÀI TẬP: XÁC ĐỊNH VÀ ĐỊNH NGHĨA VẤN ĐỀ}

    \vspace{0.3cm}

    {\normalsize
    Đề tài: SMART MONEY BOT -- Trợ lý ảo thông minh quản lý chi tiêu cá nhân
    }

    \vspace{0.25cm}

    {\small
    CSC10014 -- Computational Thinking
    $\cdot$ Project: Smart Virtual Assistant System
    $\cdot$ Nhóm 6
    }

    \vspace{0.1cm}

    {\small
    Đại diện nộp bài: Nguyễn Phạm Trung Dũng -- 25127306
    }
\end{center}

\vspace{0.3cm}

\TableStart
\begingroup
\normalsize
\begin{tabularx}{\textwidth}{
    |P{0.9cm}|P{5.6cm}|P{2.1cm}|Y|
}
\hline
\textbf{STT}
& \textbf{Họ và tên}
& \textbf{MSSV}
& \textbf{Vai trò chính}
\\
\hline
1 & Nguyễn Phạm Trung Dũng & 25127306
  & Trưởng nhóm - Làm slides
\\
\hline
2 & Đỗ Nhật Minh & 21127355
  & Phó nhóm - Viết báo cáo ý tưởng, vấn đề
\\
\hline
3 & Nguyễn Mạnh Cường & 25127024
  & Viết báo cáo ý tưởng, vấn đề
\\
\hline
4 & Hồ Thái Nguyên & 25127105
  & Viết báo cáo ý tưởng, vấn đề
\\
\hline
5 & Phan Xuân Phúc & 25127229
  & Viết báo cáo ý tưởng, vấn đề
\\
\hline
6 & Vũ Đức Duy & 25127314
  & Làm slides
\\
\hline
\end{tabularx}
\endgroup
\TableEnd

\vspace{0.3cm}

{\large\bfseries 1. Bối cảnh và phát biểu vấn đề}

\vspace{0.15cm}

{\normalsize\bfseries 1.1. Bối cảnh}

Sinh viên và người đi làm trẻ chi tiêu hằng ngày cho ăn uống,
mua sắm, đi lại và các hóa đơn định kỳ (điện, nước, internet,
điện thoại). Thông tin chi tiêu nằm rải rác ở nhiều nơi: hóa đơn
giấy, ảnh chụp, tin nhắn ngân hàng, ghi chú. Ghi chép thủ công
tốn công nên dễ bị bỏ dở.

Các ứng dụng hiện có (Sổ Thu Chi MISA, Money Lover, Spendee)
đã giải quyết một phần bài toán, nhưng qua khảo sát của nhóm,
nhiều tiện ích tự động nằm ở gói trả phí, gói miễn phí bị giới hạn
số ví/ngân sách, và việc nhập liệu tự động vẫn có thể sai số tiền,
ngày hoặc danh mục nên người dùng vẫn phải kiểm tra lại.

\vspace{0.5cm}

{\normalsize\bfseries 1.2. Phát biểu vấn đề}

Sinh viên và người đi làm trẻ khó kiểm soát chi tiêu cá nhân vì:
(1) ghi chép chi tiêu tốn công nên bỏ dở, dẫn đến không biết mình
đã tiêu bao nhiêu cho từng loại; (2) đặt ngân sách theo cảm tính
và chỉ nhận ra mình vượt mức khi đã muộn; (3) quên hạn các hóa
đơn định kỳ, dễ bị phạt hoặc gián đoạn dịch vụ.

\textbf{Câu hỏi cốt lõi:} người dùng cần gì để nắm được chi tiêu
của mình, đặt và giữ được hạn mức, không trễ hóa đơn, mà tốn ít
công sức nhất?

\newpage

\vspace{0.5cm}

{\normalsize\bfseries 1.3. Chia nhỏ thành các vấn đề con}
\vspace{0.5cm}

\TableStart
\begingroup
\normalsize
\begin{tabularx}{\textwidth}{
    |P{0.8cm}|P{4.7cm}|P{5.9cm}|Y|
}
\hline
\textbf{Mã}
& \textbf{Vấn đề con}
& \textbf{Biểu hiện}
& \textbf{Hệ quả}
\\
\hline
P1
& Ghi chép chi tiêu tốn công, dữ liệu rải rác và không đồng nhất
& Hóa đơn giấy, ảnh, tin nhắn, ghi chú; ghi tay dễ quên
& Không biết tiêu bao nhiêu cho từng loại
\\
\hline
P2
& Đặt hạn mức theo cảm tính, không có cảnh báo sớm
& Hạn mức không dựa trên chi tiêu thực tế; chỉ biết khi hết tháng
& Vượt ngân sách mà không nhận ra kịp
\\
\hline
P3
& Quên hạn hóa đơn định kỳ
& Nhiều hóa đơn, nhiều ngày đến hạn khác nhau
& Phí phạt, gián đoạn điện, nước, internet
\\
\hline
P4
& Không thể tin tuyệt đối kết quả nhập tự động
& Có thể sai số tiền, ngày, danh mục; hóa đơn mờ hoặc câu nhập thiếu thông tin
& Số liệu sai dẫn đến quyết định sai
\\
\hline
\end{tabularx}
\endgroup
\TableEnd

\vspace{0.5cm}

{\normalsize\bfseries 1.4. Phạm vi}

\textbf{Trong phạm vi:} ghi nhận chi tiêu từ hóa đơn (ảnh/PDF) và
câu tiếng Việt ngắn; phân loại vào 6 danh mục; thống kê; hạn mức
và cảnh báo; quản lý và nhắc hóa đơn định kỳ. Một người dùng,
một loại tiền (VNĐ), chủ yếu tiếng Việt.

\textbf{Ngoài phạm vi (phiên bản hiện tại):} kết nối trực tiếp tài
khoản ngân hàng/ví điện tử, thanh toán thật, chia sẻ ví nhiều
người, đa tiền tệ, tư vấn mua sắm ``đáng tiền''. Ghi nhận là
hướng mở rộng.

\vspace{0.5cm}

{\normalsize\bfseries 1.5. ``Smart'' khác gì trợ lý ảo thông thường?}

Nhóm phân biệt theo hành vi mà người dùng cảm nhận được, không
theo công nghệ:

\TableStart
\begingroup
\normalsize
\begin{tabularx}{\textwidth}{
    |P{2.6cm}|P{6.4cm}|Y|
}
\hline
\textbf{Tiêu chí}
& \textbf{Trợ lý ảo thông thường (sổ thu chi)}
& \textbf{Smart Money Bot}
\\
\hline
Nhập liệu
& Người dùng tự điền từng ô
& Hiểu câu tự nhiên và hóa đơn, điền sẵn, người dùng chỉ xác nhận
\\
\hline
Hạn mức
& Người dùng tự đặt
& Gợi ý dựa trên lịch sử chi tiêu thực tế
\\
\hline
Nhắc nhở
& Chỉ phản hồi khi được yêu cầu
& Chủ động báo trước khi sắp vượt hạn mức hoặc sắp đến hạn
\\
\hline
Khi thiếu thông tin
& Báo lỗi hoặc bỏ qua
& Nhận ra chỗ chưa chắc và hỏi lại người dùng
\\
\hline
\end{tabularx}
\endgroup
\TableEnd

Các ứng dụng hiện có đã làm được một phần các điểm trên. Điểm
nhóm hướng tới là gom cả bốn hành vi vào một luồng chủ động, và
luôn để người dùng kiểm tra trước khi lưu.

\newpage

\vspace{0.5cm}

{\large\bfseries 2. Xác định các bên liên quan (Stakeholders)}

\TableStart
\begingroup
\normalsize
\begin{tabularx}{\textwidth}{
    |P{2.8cm}|P{2.0cm}|P{3.4cm}|P{4.0cm}|Y|
}
\hline
\textbf{Bên liên quan}
& \textbf{Vai trò}
& \textbf{Nhu cầu / kỳ vọng}
& \textbf{Mối quan tâm / rủi ro}
& \textbf{Mức ảnh hưởng}
\\
\hline
Sinh viên, người đi làm trẻ (người dùng chính)
& Nhập hóa đơn, xem thống kê, đặt hạn mức, nhận nhắc nhở
& Ghi chép nhanh, ít thao tác; biết tiêu bao nhiêu cho từng loại; không trễ hạn, không vượt ngân sách
& Phân loại sai, sai số tiền; quá nhiều thông báo gây phiền; lộ dữ liệu tài chính
& Cao
\\
\hline
Người thân, bạn cùng phòng (người dùng thứ cấp)
& Cùng chia sẻ chi phí sinh hoạt chung
& Hóa đơn chung được trả đúng hạn, thông tin chi phí rõ ràng
& Bản hiện tại là đơn người dùng nên phụ thuộc người dùng chính cập nhật
& Trung bình
\\
\hline
Nhóm phát triển
& Phân tích, thiết kế, xây dựng, kiểm thử
& Sản phẩm chạy được, đúng phạm vi, đúng tiến độ 11 tuần
& Giới hạn thời gian, dữ liệu mẫu, chi phí API
& Cao
\\
\hline
Giảng viên đánh giá
& Đánh giá theo rubric
& Vận dụng đúng các bước tư duy tính toán, có thành phần thông minh, nhóm hiểu rõ sản phẩm
& Đạo văn, nội dung do AI sinh ra mà nhóm không hiểu
& Cao
\\
\hline
Nhà cung cấp dịch vụ AI/OCR/API
& Cung cấp dịch vụ đọc hóa đơn, phân loại văn bản
& Sử dụng đúng điều khoản, trong hạn mức miễn phí
& Giới hạn lượt gọi, đổi chính sách gói
& Thấp
\\
\hline
Đơn vị phát hành hóa đơn
& Nguồn tạo dữ liệu đầu vào
& Hóa đơn được đọc đúng
& Định dạng hóa đơn đa dạng làm giảm độ chính xác
& Thấp
\\
\hline
\end{tabularx}
\endgroup
\TableEnd

\vspace{0.15cm}

\textbf{Lưu ý xuyên suốt:} dữ liệu chi tiêu là dữ liệu nhạy cảm,
nên chỉ lưu tối thiểu và luôn để người dùng kiểm tra, xác nhận
trước khi lưu.

\vspace{0.5cm}
\newpage
{\large\bfseries 3. Làm rõ mục tiêu (Objectives)}

Mỗi mục tiêu là một danh sách việc con cụ thể. Mỗi việc con có
chỉ số đo, ngưỡng đạt và cách kiểm chứng (tiêu chí SMART).
Mục tiêu hoàn thành khi tất cả việc con đạt ngưỡng; có 15 việc
con, hoàn thành hết là 100\%.

\vspace{0.1cm}

\textbf{Quy ước 6 danh mục chi tiêu:}

\begin{itemize}
    \item \textbf{Ăn uống:} bữa ăn, ăn vặt, cà phê, đồ uống,
    đặt đồ ăn online.

    \item \textbf{Tiêu dùng/Mua sắm:} nhu yếu phẩm, quần áo,
    mỹ phẩm, đồ gia dụng nhỏ, thiết bị cá nhân.

    \item \textbf{Đi lại:} xăng xe, vé xe buýt, ứng dụng đặt xe,
    phí gửi xe, bảo dưỡng xe.

    \item \textbf{Hóa đơn tiện ích:} điện, nước, internet,
    cước di động định kỳ.

    \item \textbf{Giải trí:} vé xem phim, thuê bao trực tuyến,
    sự kiện, sách truyện giải trí.

    \item \textbf{Khác:} y tế, thuốc men, đồ dùng học tập đặc thù,
    các khoản phát sinh không định kỳ.
\end{itemize}

\textbf{Mục tiêu G1: Người dùng biết mình tiêu bao nhiêu cho từng
loại (giải quyết P1, P4)}

\TableStart
\begingroup
\normalsize
\begin{tabularx}{\textwidth}{
    |P{0.8cm}|P{4.0cm}|P{2.4cm}|P{4.5cm}|Y|
}
\hline
\textbf{Mã}
& \textbf{Việc con}
& \textbf{Chỉ số đo}
& \textbf{Ngưỡng đạt}
& \textbf{Cách kiểm chứng}
\\
\hline
1.1
& Ghi nhận đủ thông tin khoản chi từ hóa đơn và câu nhập tiếng Việt
& Tỷ lệ trường đúng
& Số tiền và ngày đúng $\geq 90\%$; tên cửa hàng đúng $\geq 80\%$ ($\geq 50$ hóa đơn mẫu, $\geq 20$ câu nhập mẫu)
& Đối chiếu với giá trị gán tay
\\
\hline
1.2
& Gán đúng một trong 6 danh mục
& Độ chính xác phân loại
& $\geq 85\%$ trên $\geq 50$ hóa đơn mẫu, mỗi danh mục có $\geq 5$ mẫu
& So với nhãn gán tay, lập ma trận nhầm lẫn
\\
\hline
1.3
& Khoản chi thiếu hoặc chưa chắc được đánh dấu, người dùng kiểm tra và sửa trước khi lưu
& Tỷ lệ khoản thiếu số tiền hoặc ngày bị chặn lưu
& 100\% trên $\geq 5$ test case
& Bảng test case
\\
\hline
1.4
& Thống kê chi theo danh mục, theo tháng, lịch sử $\geq 3$ tháng
& Độ khớp với tính tay
& 100\% trên 5 bộ dữ liệu mô phỏng
& Tính tay song song
\\
\hline
1.5
& Xử lý mỗi khoản chi đủ nhanh
& Thời gian xử lý
& $\leq 10$ giây mỗi hóa đơn
& Đo thời gian chạy
\\
\hline
\end{tabularx}
\TableEnd

\newpage

\normalsize\textbf{Mục tiêu G2: Đặt hạn mức dựa trên thực tế và được báo khi
sắp vượt (giải quyết P2)}

\TableStart
\begingroup
\normalsize
\begin{tabularx}{\textwidth}{
    |P{0.8cm}|P{4.0cm}|P{2.7cm}|P{4.4cm}|Y|
}
\hline
\textbf{Mã}
& \textbf{Việc con}
& \textbf{Chỉ số đo}
& \textbf{Ngưỡng đạt}
& \textbf{Cách kiểm chứng}
\\
\hline
2.1
& Gợi ý hạn mức cho từng danh mục từ $\geq 3$ tháng dữ liệu
& Độ khớp với công thức tham chiếu của nhóm; thời gian phản hồi
& Khớp 100\% trên 5 bộ dữ liệu; $\leq 5$ giây
& Tính tay song song, đo thời gian
\\
\hline
\textbf{Mã}
& \textbf{Việc con}
& \textbf{Chỉ số đo}
& \textbf{Ngưỡng đạt}
& \textbf{Cách kiểm chứng}
\\
\hline
2.2
& Người dùng nhập ngân sách tổng và chỉnh hạn mức; được báo nếu tổng hạn mức vượt ngân sách tổng
& Số test case đúng
& 100\% trên $\geq 5$ test case
& Bảng test case
\\
\hline
2.3
& Báo khi chi tiêu đạt 80\% và khi đạt hoặc vượt 100\% hạn mức của một danh mục
& Tỷ lệ cảnh báo đúng (không bỏ sót, không báo nhầm)
& 100\% trên $\geq 10$ test case, gồm trường hợp một khoản chi nhảy qua cả hai mốc
& Bảng kịch bản: đầu vào, cảnh báo mong đợi, kết quả thực tế
\\
\hline
2.4
& Người dùng đổi được mốc cảnh báo
& Số test case đúng
& 100\% trên $\geq 3$ test case
& Bảng test case
\\
\hline
\end{tabularx}
\endgroup
\TableEnd

\vspace{0.25cm}

\textbf{Mục tiêu G3: Người dùng không trễ hóa đơn định kỳ
(giải quyết P3)}

\vspace{0.2cm}

\TableStart
\begingroup
\normalsize
\begin{tabularx}{\textwidth}{
    |P{0.8cm}|P{4.0cm}|P{2.5cm}|P{4.5cm}|Y|
}
\hline
\textbf{Mã}
& \textbf{Việc con}
& \textbf{Chỉ số đo}
& \textbf{Ngưỡng đạt}
& \textbf{Cách kiểm chứng}
\\
\hline
3.1
& Lưu hóa đơn định kỳ: tên dịch vụ, số tiền dự kiến, ngày đến hạn, chu kỳ lặp
& Số hóa đơn nhập đủ trường
& 100\% trên $\geq 5$ hóa đơn mẫu
& Kiểm tra dữ liệu lưu
\\
\hline
3.2
& Nhắc trước 3 ngày và trước 1 ngày
& Tỷ lệ nhắc đúng thời điểm
& 100\% trên $\geq 10$ test case, gồm ngày giáp ranh tháng và hạn ngày 31 ở tháng ít ngày
& Mô phỏng đổi ngày hệ thống, kiểm tra thông báo
\\
\hline
3.3
& Hiển thị đúng trạng thái Đã thanh toán / Sắp đến hạn / Quá hạn; hóa đơn đã thanh toán không bị nhắc nữa
& Số test case đúng
& 100\% trên $\geq 6$ test case
& Bảng test case
\\
\hline
\end{tabularx}
\endgroup
\TableEnd

\newpage

\textbf{Mục tiêu G4: Giảm công sức ghi chép và có sản phẩm kiểm
chứng được (giải quyết vấn đề cốt lõi)}

\TableStart
\begingroup
\normalsize
\begin{tabularx}{\textwidth}{
    |P{0.8cm}|P{4.0cm}|P{2.7cm}|P{4.5cm}|Y|
}
\hline
\textbf{Mã}
& \textbf{Việc con}
& \textbf{Chỉ số đo}
& \textbf{Ngưỡng đạt}
& \textbf{Cách kiểm chứng}
\\
\hline
4.1
& Ghi một khoản chi với ít thao tác
& Số thao tác từ lúc đưa hóa đơn/câu nhập đến lúc lưu
& $\leq 3$ thao tác khi dữ liệu được đọc đúng
& Thử trực tiếp, đếm thao tác
\\
\hline
4.2
& Tiết kiệm thời gian so với ghi tay
& Thời gian ghi một khoản chi
& Giảm $\geq 50\%$ so với ghi tay ($\geq 10$ khoản, do các thành viên nhóm thử)
& Bấm giờ hai cách
\\
\hline
4.3
& Có bản chạy được cả 3 nhóm chức năng và đã kiểm thử
& Số chức năng chạy được; số test case
& 3/3 chức năng trước cuối tuần 7; $\geq 8$ test case, trong đó $\geq 5$ trường hợp biên
& Demo trực tiếp, bảng test case trong báo cáo
\\
\hline
\end{tabularx}
\endgroup
\TableEnd

\vspace{0.15cm}

\textbf{Ghi chú:} các ngưỡng là mục tiêu của nhóm cho phiên bản
mô phỏng, dùng dữ liệu do nhóm chuẩn bị; có thể điều chỉnh sau
khảo sát dữ liệu và được ghi lại trong phần Testing \& Reflection.

\vspace{0.3cm}

{\large\bfseries 4. Xác định đầu vào và đầu ra
(Inputs / Expected Outputs)}

\vspace{0.15cm}

{\normalsize\bfseries 4.1. Đầu vào}

\vspace{0.5cm}

\TableStart
\begingroup
\normalsize
\begin{tabularx}{\textwidth}{
    |P{0.8cm}|P{3.0cm}|P{6.1cm}|Y|
}
\hline
\textbf{Mã}
& \textbf{Đầu vào}
& \textbf{Mô tả / định dạng}
& \textbf{Ví dụ}
\\
\hline
I1
& Hóa đơn dạng ảnh/PDF
& File .jpg, .png, .pdf của hóa đơn bán lẻ hoặc hóa đơn dịch vụ
& Hóa đơn quán cà phê, siêu thị
\\
\hline
I2
& Giao dịch nhập bằng văn bản tiếng Việt
& Câu ngắn tự nhiên có số tiền và nội dung chi
& ``ăn trưa 45k'', ``mua sạc điện thoại 150k''
\\
\hline
I3
& Lịch sử chi tiêu
& Bảng gồm ngày, số tiền, danh mục, ghi chú; tối thiểu 3 tháng (dữ liệu mô phỏng)
& Bảng CSV 90 ngày
\\
\hline
\end{tabularx}
\endgroup
\TableEnd

\TableStart
\begingroup
\normalsize
\begin{tabularx}{\textwidth}{
    |P{0.8cm}|P{3.0cm}|P{6.1cm}|Y|
}
\hline
\textbf{Mã}
& \textbf{Đầu vào}
& \textbf{Mô tả / định dạng}
& \textbf{Ví dụ}
\\
\hline
I4
& Ngân sách tổng mỗi tháng
& Một số tiền (VNĐ) do người dùng nhập
& 6.000.000 đ/tháng
\\
\hline
I5
& Hóa đơn định kỳ
& Tên hóa đơn, số tiền, ngày đến hạn, chu kỳ lặp
& Tiền điện, 850.000 đ, hạn ngày 20 hằng tháng
\\
\hline
I6
& Thiết lập của người dùng
& Mốc cảnh báo (mặc định 80\%, 100\%), số ngày nhắc trước (mặc định 3 và 1), hạn mức chỉnh tay
& Đổi mốc cảnh báo thành 70\%
\\
\hline
I7
& Xác nhận của người dùng
& Xác nhận hoặc sửa khoản chi trước khi lưu; đánh dấu hóa đơn đã thanh toán
& Sửa số tiền, bấm ``Đã thanh toán''
\\
\hline
\end{tabularx}
\endgroup
\TableEnd

\vspace{0.5cm}

{\normalsize\bfseries 4.2. Đầu ra mong đợi}

\TableStart
\begingroup
\normalsize
\begin{tabularx}{\textwidth}{
    |P{1cm}|P{2.5cm}|P{7cm}|P{4.5cm}|Y|
}
\hline
\textbf{Mã}
& \textbf{Đầu ra}
& \textbf{Mô tả}
& \textbf{Ví dụ}
& \textbf{Phục vụ}
\\
\hline
Ou1
& Bản ghi giao dịch
& Cửa hàng, ngày, số tiền, danh mục; trường chưa chắc được đánh dấu để người dùng xác nhận hoặc sửa trước khi lưu
& Highlands Coffee, 13/10/2026, 65.000 đ, Ăn uống
& G1
\\
\hline
Ou2
& Thống kê chi tiêu
& Bảng và biểu đồ tổng chi theo danh mục, theo tháng
& Cơ cấu chi tiêu tháng 10
& G1
\\
\hline
Ou3
& Hạn mức đề xuất
& Hạn mức từng danh mục, người dùng chỉnh được; báo nếu tổng vượt ngân sách tổng
& Ăn uống: 2.000.000 đ/tháng
& G2
\\
\hline
Ou4
& Cảnh báo ngân sách
& Thông báo khi đạt 80\% hoặc đạt/vượt 100\% hạn mức
& ``Ăn uống đã dùng 80\% hạn mức tháng này''
& G2
\\
\hline
Ou5
& Nhắc hạn hóa đơn
& Danh sách hóa đơn sắp đến hạn và quá hạn, kèm số ngày còn lại
& ``Hóa đơn điện 850.000 đ đến hạn sau 3 ngày''
& G3
\\
\hline
Ou6
& Trạng thái hóa đơn
& Đã thanh toán / Sắp đến hạn / Quá hạn, quá hạn được cảnh báo nổi bật
& Hóa đơn nước: Quá hạn 2 ngày
& G3
\\
\hline
\end{tabularx}
\endgroup
\TableEnd

\newpage

\vspace{0.5cm}

{\normalsize\bfseries 4.3. Ví dụ minh họa}

\textbf{Trường hợp chuẩn.} Người dùng đưa hóa đơn
``Highlands Coffee -- 65.000 đ -- 13/10/2026''. Hạn mức Ăn uống
tháng 10 là 2.000.000 đ, đã chi 1.540.000 đ. Sau khi người dùng
xác nhận và lưu, tổng Ăn uống là 1.605.000 đ, tương đương
80,25\% hạn mức, nên người dùng nhận cảnh báo
``Ăn uống đã đạt 80\% hạn mức''. Hóa đơn điện 850.000 đ hạn
20/10 được nhắc vào 17/10 và 19/10.

\textbf{Trường hợp ngoại lệ.} Hóa đơn bị rách hoặc mờ chữ số
tiền, hoặc câu nhập thiếu như ``mua đồ bách hóa xanh''. Thông
tin đã rõ (cửa hàng Bách Hóa Xanh, danh mục Tiêu dùng/Mua sắm)
được giữ lại; số tiền chưa chắc được đánh dấu và người dùng
phải điền trước khi lưu.

\vspace{0.5cm}

{\normalsize\bfseries 5. Nêu các ràng buộc (Constraints)}

\TableStart
\begingroup
\normalsize
\begin{tabularx}{\textwidth}{
    |P{3.5cm}|P{9cm}|Y|
}
\hline
\textbf{Loại}
& \textbf{Nội dung ràng buộc}
& \textbf{Ảnh hưởng đến bài toán}
\\
\hline
Thời gian
& Dự án 11 tuần; làm bản chạy thử tuần 6--7; kiểm thử tuần 8--9; báo cáo và bảo vệ tuần 10--11
& Chỉ chọn 3 nhóm chức năng cốt lõi
\\
\hline
Nhân lực
& Nhóm 6 sinh viên; mọi thành viên tham gia các giai đoạn (phân tích, thiết kế, lập trình, báo cáo) và ghi logbook hằng tuần; mỗi người có vai trò chính nhưng đều phải hiểu toàn bộ nội dung
& Phân công đồng đều, mỗi người phụ trách một phần
\\
\hline
Kỹ thuật
& Mô phỏng ở mức prototype bằng Python, Streamlit, Pandas và API AI/OCR; chạy cục bộ; không kết nối ngân hàng hoặc ví điện tử thật
& Giao dịch đến từ hóa đơn và nhập tay, không đồng bộ tự động
\\
\hline
Độ chính xác
& Kết quả đọc hóa đơn và phân loại phụ thuộc chất lượng ảnh và kiểu hóa đơn; có thể sai số tiền, ngày, danh mục
& Bắt buộc để người dùng kiểm tra và xác nhận trước khi lưu
\\
\hline
Dữ liệu
& Dùng dữ liệu mô phỏng và hóa đơn mẫu do nhóm chuẩn bị ($\geq 50$ hóa đơn, $\geq 3$ tháng lịch sử); không dùng dữ liệu tài chính thật của người khác
& Kết quả đánh giá phản ánh trên bộ dữ liệu mẫu có kiểm soát
\\
\hline
Chi phí
& Chỉ dùng gói miễn phí hoặc credit giới hạn của API
& Giới hạn số lần gọi API
\\
\hline
\end{tabularx}
\endgroup
\TableEnd

\newpage

% ============================================================
% TRANG 10
% ============================================================

\TableStart
\begingroup
\normalsize
\begin{tabularx}{\textwidth}{
    |P{3.2cm}|P{9cm}|Y|
}
\hline
\textbf{Loại}
& \textbf{Nội dung ràng buộc}
& \textbf{Ảnh hưởng đến bài toán}
\\
\hline
Phạm vi
& Hóa đơn chủ yếu tiếng Việt; một loại tiền (VNĐ); một người dùng; 6 danh mục
& Không xử lý đa ngôn ngữ, đa tiền tệ, chia sẻ ví
\\
\hline
Quyền riêng tư
& Không lưu số thẻ, số tài khoản hoặc thông tin định danh nhạy cảm; chỉ dùng hóa đơn mẫu không chứa thông tin cá nhân thật khi gọi dịch vụ AI bên ngoài
& Lưu tối thiểu: ngày, số tiền, danh mục, ghi chú
\\
\hline
Học thuật
& Được dùng ChatGPT, Claude, Gemini để lên ý tưởng và tham khảo, nhưng nội dung và mã phải được nhóm hiểu, chỉnh sửa và viết lại; đạo văn bị 0 điểm
& Mọi thành viên phải giải thích được khi bảo vệ
\\
\hline
\end{tabularx}
\endgroup
\TableEnd

\vspace{0.5cm}

{\normalsize\bfseries 6. Tóm tắt}

Vấn đề: người dùng trẻ khó kiểm soát chi tiêu vì ghi chép tốn
công, đặt hạn mức cảm tính và quên hạn hóa đơn. Vấn đề được
chia thành 4 vấn đề con (P1--P4), đối ứng với 4 mục tiêu
(G1--G4) gồm 15 việc con đo được; hoàn thành hết là đạt 100\%.
Hệ thống được phân biệt với trợ lý ảo thông thường qua 4 hành
vi thông minh (mục 1.5), có 7 nhóm đầu vào (I1--I7), 6 nhóm
đầu ra (Ou1--Ou6) và chịu các ràng buộc về thời gian, nhân lực,
kỹ thuật, độ chính xác, dữ liệu, chi phí, phạm vi, quyền riêng tư
và học thuật.

\end{document}